% continuum_limits.tex
% Section 8 -- discrete-to-continuum limit operators.  The Hilbert-Sixth
% on-ramp.  Lattice-spacing parameter epsilon, epsilon -> 0 limit,
% convergence modes, consistency conditions tying discrete to
% continuum operators.

\label{sec:continuum_limits}

\placeholder{Section body -- write during Phase 4.  This is the
load-bearing section for the planned Hilbert-Sixth capstone.  Develop
the discrete-to-continuum limit machinery: parameter, topology,
convergence modes, consistency theorems.  The most mathematical
care in the paper goes here.}

\subsection{The Limit Parameter}
\label{subsec:limit_parameter}

\placeholder{Introduce the lattice-spacing parameter $\epsilon > 0$;
the rescaled lattice $\epsilon \Lambda_d \subset \mathbb{R}^d$; the
$\epsilon \to 0$ limit as the discrete-to-continuum passage.}

\subsection{Convergence Modes}
\label{subsec:convergence_modes}

\placeholder{Define what it means for a sequence of discrete
operators $T_\epsilon \colon \ell^p(\epsilon \Lambda_d) \to
\ell^p(\epsilon \Lambda_d)$ to converge to a continuum operator
$T \colon L^p(\mathbb{R}^d) \to L^p(\mathbb{R}^d)$.  Three modes:
strong, weak, and distributional.  Relations among them.}

\subsection{Consistency Conditions}
\label{subsec:consistency_conditions}

\placeholder{Conditions tying discrete operators to their continuum
limits.  The discrete exterior derivative $\mathrm{d}$
(Section~\ref{sec:discrete_diff_geometry}) converges to the
continuum $\mathrm{d}$ in $L^p$ for appropriate function classes;
the discrete Hodge Laplacian $\Delta_{\Lambda_d}$ converges to the
continuum Laplacian; \ldots.  These are the consistency theorems
the Hilbert-Sixth capstone will cite.}

\subsection{Function-Space Convergence}
\label{subsec:function_space_convergence}

\placeholder{Convergence of the function spaces themselves:
$\ell^p(\epsilon \Lambda_d) \to L^p(\mathbb{R}^d)$ in an
appropriate sense; Sobolev-embedding constants stabilising in the
limit.}

\subsection{Why This is the Hilbert-Sixth On-Ramp}
\label{subsec:hilbert_sixth_onramp}

\placeholder{Hilbert's 1900 problem specifically singled out
Boltzmann's kinetic theory and the rigorous derivation of fluid
equations as the centerpiece of an axiomatised physics.  Both
require a discrete-to-continuum limit treated with full rigour.
This section provides the apparatus the capstone synthesises into
the axiomatisation claim.  Cross-reference: the motivating examples
in Section~\ref{sec:motivating_examples} show what the apparatus is
used \emph{for}; the dimensional-selection corollary in
Section~\ref{sec:dimensional_selection} is the other principal
consumer.}
